% Generated draft: requires mathematical and source review.
\documentclass[11pt]{article}
\usepackage[margin=25mm]{geometry}
\usepackage{fontspec}
\setmainfont{Latin Modern Roman}
\usepackage{amsmath,amssymb,mathtools,unicode-math}
\setmathfont{Latin Modern Math}
\usepackage{enumitem,needspace,xurl,hyperref}
\setlength{\parindent}{0pt}
\setlength{\parskip}{5pt}
\setlength{\emergencystretch}{3em}
\newcommand{\N}{\mathbb{N}}
\newcommand{\Z}{\mathbb{Z}}
\newcommand{\Q}{\mathbb{Q}}
\newcommand{\R}{\mathbb{R}}
\newcommand{\C}{\mathbb{C}}
\newcommand{\F}{\mathbb{F}}
\newcommand{\A}{\mathbb{A}}
\newcommand{\PP}{\mathbb{P}}
\newcommand{\E}{\mathbb{E}}
\newcommand{\eps}{\varepsilon}
\newcommand{\dd}{\,\mathrm{d}}
\begin{document}
\section*{UnsolvedMath ID: 801105}
\textbf{Source problem code:} catalogue entry
\subsection{Definitions and mathematical statement}
For a parameterized interpretation of the displayed construction, let \(I\) be an index set, let \(\mathcal U\) be an ultrafilter on \(I\), and let \((n_i)_{i\in I}\) be positive integers. These parameters make the construction explicit; the supplied record does not specify their hypotheses or quantification. Write \(\mathbb S_{n_i}\) for the group of permutations of \(\{1,\ldots,n_i\}\), equipped with the normalized Hamming distance
\[
d_{H,i}(\sigma,\rho)=\frac{1}{n_i}\bigl|\{a:\sigma(a)\ne\rho(a)\}\bigr|.
\]
Its metric ultraproduct is
\[
\mathbb S^{\mathcal U}=\left(\prod_{i\in I}\mathbb S_{n_i}\right)\Big/\left\{(\sigma_i):\lim_{i\to\mathcal U}d_{H,i}(\sigma_i,e_i)=0\right\},
\]
where \(e_i\) is the identity permutation.

The standard tracial interpretation of the source's ambient matrix ultraproduct is
\[
\mathbb M^{\mathcal U}=\frac{\{(x_i):x_i\in M_{n_i}(\C),\ \sup_i\|x_i\|<\infty\}}{\{(x_i):\sup_i\|x_i\|<\infty,\ \lim_{i\to\mathcal U}\|x_i\|_{2,i}=0\}},
\qquad
\|x\|_{2,i}^{2}=\operatorname{tr}_{n_i}(x^*x),
\]
with \(\operatorname{tr}_{n_i}=n_i^{-1}\operatorname{Tr}\) and trace \(\tau([(x_i)])=\lim_{i\to\mathcal U}\operatorname{tr}_{n_i}(x_i)\). The supplied record does not itself define \(\mathbb M^{\mathcal U}\); this is an explicit conventional interpretation.

Sending each permutation to its permutation matrix gives the source's embedding \(\mathbb S^{\mathcal U}\subset\mathbb M^{\mathcal U}\). Indeed, for permutation matrices,
\[
\|P_{\sigma}-P_{\rho}\|_{2,i}^{2}=2d_{H,i}(\sigma,\rho),
\qquad
\tau([(P_{\sigma_i})])=1-\lim_{i\to\mathcal U}d_{H,i}(\sigma_i,e_i).
\]
Let \(A=W^*(\mathbb S^{\mathcal U})\) be the von Neumann subalgebra of \(\mathbb M^{\mathcal U}\) generated by this image.

\textbf{Source question.} Does \(A\) contain a copy of a matrix ultraproduct?

Here a tracial matrix ultraproduct has the form \(\prod_{j\to\mathcal V}(M_{k_j}(\C),\operatorname{tr}_{k_j})\), with the same bounded-sequence and vanishing-\(2\)-norm quotient construction. A fully specified embedding question requires the intended restrictions on \((k_j)\) and \(\mathcal V\), and the intended meaning of a copy. The supplied record does not say whether matrix dimensions must tend to infinity, whether the target must use the original dimensions or ultrafilter, or whether the embedding must be unital, normal, or trace preserving. It also does not explicitly impose \(n_i\to\mathcal U\infty\) or state the quantifier over the initial ultrafilter and dimensions. These qualifications remain unresolved; allowing arbitrary bounded matrix dimensions could substantially change the question.
\subsection{Short English statement}
Elayavalli's Problem 6 asks whether the von Neumann algebra generated by the symmetric-group metric ultraproduct inside the ambient tracial matrix ultraproduct contains a matrix ultraproduct. The supplied record labels the problem open. The intended dimension growth, ultrafilter quantifiers, and embedding requirements need clarification before a unique precise formulation can be fixed.
\subsection{Sources}
\begin{itemize}
\item Supplied local UnsolvedMath registry record, permanent ID \(801105\), exact source problem code \texttt{catalogue entry}, titled \emph{Ultraproduct II\(_1\) Factors --- Problem 6}. The record attributes the problem to Elayavalli and gives proposed year 2025. Its status is supplied as \texttt{open}, not independently verified here.
\item Supplied source attribution: Elayavalli, \emph{50 Open Problems: Ultraproduct II\(_1\) Factors} (2025), Problem 6, \url{https://arxiv.org/abs/2511.20377}. The background records access on 2026-10-16 and extraction by \texttt{arxiv\_latex}. These are supplied provenance details; no independent source visit or verification is asserted.
\end{itemize}
\end{document}
